\documentclass[10pt,a4paper]{amsart}
\usepackage[margin=19mm]{geometry}
\usepackage{amsmath,amssymb,mathtools}
\usepackage[scr=rsfs,cal=euler]{mathalfa}
\usepackage{xfrac}
\usepackage[unicode,hidelinks,bookmarks=false]{hyperref}
\newcommand{\ie}{i.e.\ }
\newcommand{\cf}{cf.\ }
\newcommand{\resp}{respectively\ }
\newcommand{\Subsection}{\subsection}
\providecommand{\nobreakdash}{\nobreak\hbox{-}}
\setlength{\emergencystretch}{2em}
\multlinegap0pt
\usepackage{amssymb}
\usepackage{float}
\usepackage{dsfont}
\usepackage{enumerate}
\usepackage{algorithm}\usepackage{algpseudocode}
\usepackage{tikz}
\newtheorem{assumption}{Assumption}
\newtheorem{lemma}{Lemma}
\usepackage{cleveref}
\crefname{assumption}{Assumption}{Assumptions}
\crefname{lemma}{Lemma}{Lemmas}
\def\C{\mathbb{C}}
\def\N{\mathbb{N}}
\renewcommand{\Re}{{\rm Re}}
\def\dist{{\rm dist}}
\newcommand{\inner}[1]{{\left\langle#1\right\rangle}}
\newcommand{\norm}[1]{{\left\|#1\right\|}}
\newcommand{\sampset}{X}
\renewcommand{\subset}{\subseteq}
\def\wherespace{\quad\text{where}\quad}
\title{Multidimensional Gradient-MUSIC: A Global Nonconvex Optimization Framework for Optimal Resolution}
\author{Albert Fannjiang, Weilin Li}
\date{Source: arXiv:2603.27379v1 (CC BY 4.0)}
\begin{document}
\maketitle

\noindent\emph{Source and licence.} Multidimensional Gradient-MUSIC: A Global Nonconvex Optimization Framework for Optimal Resolution. Version arXiv:2603.27379v1 (CC BY 4.0), distributed by arXiv under the Creative Commons Attribution 4.0 International licence, http://creativecommons.org/licenses/by/4.0/, as recorded in the arXiv licence metadata for this exact version and shown on the abstract page https://arxiv.org/abs/2603.27379v1. Full licence notice: \texttt{licenses/arxiv-2603.27379-CC-BY-4.0.txt}.\par\smallskip

\noindent\textit{Editorial scope.} Counted unit: the article's lemma lem:approx2 (source0.tex lines 2136--2143). The article's complete original proof of it is delivered with it (source0.tex lines 2145--2225, one \texttt{proof} environment of 81 lines). No mathematics is written, repaired or re-derived: the statement and the proof are the article's own lines, in the article's own notation, and no retained line is substituted or re-flowed.  Retained uncounted as the source-local context the proof needs: lines 765--777; lines 884--898; lines 1206--1209; lines 2122--2129.  Nothing was omitted from any retained span, so the package carries no omission record.  No retained line contains a citation, so the wrapper carries no bibliography and no bibliography entry.  No retained line contains a figure environment, an image-inclusion command or drawing code, so no image is delivered and the package declares no asset dependency.  Editorial formatting only, in addition: an AMS \texttt{amsart} preamble that restates the article's own declarations and macros used by the retained lines, namely the environments \texttt{assumption}, \texttt{lemma} on separate counters, together with the standard packages those lines need.  The retained environments print in this document as Assumption 1, Lemma 1 in order of appearance, whereas the article prints the same items with the numbering of its own full document.  The complete native source of the article as distributed in the arXiv e-print for this exact version is bundled unchanged as \texttt{sources/197342/source0.tex}.
\begin{assumption}[Symmetry and regularity]
\label{assump:symmetry}
Assume that $X=-X$, that
\[
\nu(A)=\nu(-A)
\qquad\text{for every $\nu$-measurable set } A\subset X,
\]
and that for every multi-index $\alpha\in\N^d$ with $|\alpha|\le 4$,
\begin{equation}
\label{eq:moment}
\int_X |x^\alpha|\,d\nu(x)<\infty.
\end{equation}
\end{assumption}

\begin{align}
E_0(\vartheta)
&:=
\sup_{\omega\in\Omega}
\sum_{\ell:\,|\omega-\theta_\ell|\ge \Delta/2}
|K(\omega-\theta_\ell)|^2,
\label{eq:E0def}
\\
E_1(\vartheta)
&:=
\sup_{\omega\in\Omega}
\sum_{\ell:\,|\omega-\theta_\ell|\ge \Delta/2}
|\nabla K(\omega-\theta_\ell)|^2,
\label{eq:E1def}
\end{align}

\begin{align}
	T_\vartheta\colon \C^s\to L^2_\nu(X), &\wherespace T_\vartheta v(x) := \sum_{\ell=1}^s v_\ell e^{2\pi i \theta_\ell\cdot x} \in L^2_\nu(X), \label{eq:synthesis1} \\
	T_\vartheta'\colon \C^s\to L^2_{\nu'}(X'), &\wherespace (T_\vartheta') v(x) := \sum_{\ell=1}^s v_\ell e^{2\pi i \theta_\ell\cdot x} \in L^2_{\nu'}(X'). \label{eq:synthesis2}
\end{align}

	\begin{align}
		q_W(\omega)
		&=1- \inner{\phi_\omega, P_W \phi_\omega}_{L^2_\nu}, \label{eq:q0} \\
		\partial_j q_W(\omega)
		&=-2\Re \inner{\partial_{\omega_j} \phi_\omega, P_{W} \phi_\omega}_{L^2_\nu}, \label{eq:q1} \\
		\partial_j \partial_k q_{W}(\omega)
		&=-2\Re \inner{\partial_{\omega_j} \partial_{\omega_k} \phi_\omega, P_{W} \phi_\omega}_{L^2_\nu} - 2 \Re \inner{\partial_{\omega_k} \phi_\omega, P_W \partial_{\omega_j}\phi_\omega}_{L^2_\nu}. \label{eq:q2}
	\end{align}

	\begin{lemma}
	\label{lem:approx2}
	Suppose \cref{assump:symmetry} holds. For any $\vartheta:=\{\theta_\ell\}_{\ell=1}^s$, $\tau\leq \Delta(\vartheta)/2$, and $\omega\in \Omega$ such that $\dist(\omega,\vartheta)\leq \tau$, we have
	\begin{align*}
		\|\nabla^2 q(\omega)-2 \Psi \|_F
		\leq \sqrt{2 \tau \|\nabla K\|_{L^\infty(B_\tau)} \Delta^2 K(0)} + \frac{1}{\lambda_{\min}(K_\vartheta)} \left(\|\nabla K\|_{L^\infty(B_\tau)}^2 + E_1(\vartheta) \right).
	\end{align*}
	\end{lemma}

	\begin{proof}
	Recall the shorthand notation $\Psi=-\nabla^2 K(0)$, which is a real symmetric matrix. Starting with formula \eqref{eq:q2} and doing some  further manipulations,
	\begin{equation*}
		\begin{split}
			\frac 12 \partial_j \partial_k q(\omega)
			&=-\Re \inner{\partial_{\omega_j} \partial_{\omega_k} \phi_\omega, \phi_\omega}_{L^2_\nu} + \Re \inner{\partial_{\omega_j} \partial_{\omega_k} \phi_\omega, (I-P_U) \phi_\omega}_{L^2_\nu} -\Re \inner{\partial_{\omega_j} \phi_\omega, P_U \partial_{\omega_k} \phi_\omega}_{L^2_\nu} \\ 
			&=\Psi_{j,k} + \Re \inner{\partial_{\omega_j} \partial_{\omega_k} \phi_\omega, (I-P_U) \phi_\omega}_{L^2_\nu} - \Re \inner{\partial_{\omega_j} \phi_\omega, P_U \partial_{\omega_k}\phi_\omega}_{L^2_\nu}.
		\end{split}
	\end{equation*}
	The first term is dominant, while the second and third terms are error quantities. Define the matrices $A(\omega)$ and $B(\omega)$ such that 
	\begin{align*}
		A_{j,k}(\omega)
		&=\Re \inner{\partial_{\omega_j} \partial_{\omega_k} \phi_\omega, (I-P_U) \phi_\omega}_{L^2_\nu} , \\
		B_{j,k}(\omega)
		&=- \Re \inner{\partial_{\omega_j} \phi_\omega, P_U \partial_{\omega_k}\phi_\omega}_{L^2_\nu}. 
	\end{align*}
	Hence, we have derived the formula
	\begin{equation}\label{eq:hessianq}
		\nabla^2 q (\omega) 
		= 2\Psi + 2A(\omega) + 2B(\omega). 
	\end{equation}		
	From here onward, fix any $\ell\in \{1,\dots,s\}$ and assume that $\omega\in B_\tau(\theta_\ell)$. 
	
	For the first matrix $A(\omega)$, notice that $(I-P_U)\phi_{\theta_\ell}=0$ because $P_U\phi_{\theta_\ell} = \phi_{\theta_\ell}$. For each $j,k\in \{1,\dots,d\}$,  
	\begin{align*}
		|A_{j,k}(\omega)|
		&=\left| \Re\inner{\partial_{\omega_j} \partial_{\omega_k} \phi_\omega,  (I-P_U) (\phi_\omega-\phi_{\theta_\ell})}_{L^2_\nu} \right| \\
		&\leq \norm{\partial_{\omega_j} \partial_{\omega_k} \phi_\omega}_{L^2_\nu}  \norm{\phi_\omega-\phi_{\theta_\ell}}_{L^2_\nu}.
	\end{align*}
	A calculation together with the mean value theorem and that $|\omega-\theta_\ell|\leq \tau$ shows that
	\begin{align*} 
		\norm{\phi_{\omega}-\phi_{\theta_\ell}}_{L^2_\nu}^2 
		&= 2K(0)-2K(\omega-\theta_\ell)
		\leq 2 \tau \|\nabla K\|_{L^\infty(B_\tau)}, \\
		\norm{\partial_{\omega_j} \partial_{\omega_k} \phi_\omega}_{L^2_\nu}^2
		&=\partial_j^2 \partial_k^2 K(0).  
	\end{align*}
	Using the above inequalities, we can conclude that 
	\begin{equation}
		\|A(\omega)\|_F^2
		\leq 2 \tau \|\nabla K\|_{L^\infty(B_\tau)} \sum_{j,k}  \partial_j^2 \partial_k^2 K(0)
		= 2 \tau \|\nabla K\|_{L^\infty(B_\tau)} \Delta^2 K(0).
		\label{eq:hessianq2}
	\end{equation}
	
	For the second matrix $B(\omega)$, recall the synthesis operator $T_\vartheta$ defined in \eqref{eq:synthesis1}. For convenience, let $T=T_\vartheta/\sqrt{\nu(\sampset)}$ be a normalized version and so
	$
	P_U=T (T^*T)^{-1} T^*$. A calculation yields
	\begin{align*}
		|B_{j,k}(\omega)|
		&\leq \left|\inner{\partial_{\omega_j} \phi_\omega, T (T^*T)^{-1} T^*\partial_{\omega_k}\phi_\omega}_{L^2_\nu}\right| \\
		&=\left|\inner{T^*\partial_{\omega_j} \phi_\omega, (T^*T)^{-1} (T^*\partial_{\omega_k} \phi_\omega)}_{\C^s}\right| \\
		&\leq \norm{ (T^*T)^{-1}} |T^*\partial_{\omega_j} \phi_\omega| |T^*\partial_{\omega_k} \phi_\omega|.
	\end{align*}
	Note that $T^*T = K_\vartheta$ and $T^* \partial_{\omega_j} \phi_\omega
	= \begin{bmatrix}
		\partial_j K(\omega-\theta_\ell)
	\end{bmatrix}_{\ell=1,\dots,s}
	$. Recall that $|\omega-\omega_\ell|\leq \tau$ and the definition of $E_1(\vartheta)$ defined in \eqref{eq:E1def}. Then,
	\begin{equation}
		\begin{split}
			\|B(\omega)\|_F
			&\leq \norm{ (T^*T)^{-1}} \left(\sum_{j=1}^d |T^*\partial_{\omega_j} \phi_\omega|^2 \sum_{k=1}^d |T^*\partial_{\omega_k} \phi_\omega|^2\right)^{1/2} \\
			&\leq \frac{1}{\lambda_{\min}(K_\vartheta)} \sum_{j=1}^s |\nabla K(\omega-\theta_j)|^2.
		\end{split}
		\label{eq:hessianq4}
	\end{equation}	
	We next control the summation in \eqref{eq:hessianq4} by breaking it into two terms. For all $j\ne \ell$, recalling that $|\omega-\theta_\ell|\leq \tau$, we have 
	$
	|\omega-\theta_j|
	\geq |\theta_j-\theta_\ell| - |\omega-\theta_\ell|
	\geq \Delta-\tau\geq \Delta/2. 
	$
	By definition of $E_1(\vartheta,\tau)$, this now implies 
	\begin{equation}
		\sum_{j=1}^s |\nabla K(\omega-\theta_j)|^2
		=|\nabla K(\omega-\theta_\ell)|^2 + \sum_{j\ne\ell} |\nabla K(\omega-\theta_j)|^2
		\leq \|\nabla K\|_{L^\infty(B_\tau)}^2 + E_1(\vartheta). \label{eq:hessianq5}
	\end{equation}
	Combining \eqref{eq:hessianq}, \eqref{eq:hessianq2}, \eqref{eq:hessianq4}, and \eqref{eq:hessianq5} completes the proof.
	\end{proof}

\end{document}
